\section{抛体运动的方程与轨迹}
\subsection{共同时间连接两个分运动}
忽略空气阻力，在范围不大、重力加速度可视为恒矢量的区域内，抛体满足 $\bm r=\bm r_0+\bm v_0t+\bm gt^2/2$。取水平向右、竖直向上为正，若初速度与水平线成 $\alpha$，则
\[x=x_0+v_0\cos\alpha\,t,\qquad y=y_0+v_0\sin\alpha\,t-\frac12gt^2.\]
两个方程中的时间必须相同。水平运动不能单独选一个飞行时间；是否落地、碰壁或碰到斜面，由相应边界条件确定。若 $v_0\cos\alpha\ne0$，消去时间可得抛物线轨迹；竖直抛射则应直接使用时间方程，不能除以零的水平速度。

\fig{k4deep:time}{同一组时刻在水平、竖直方向产生对应的位置}{\begin{tikzpicture}
\draw[axis](0,0)--(6,0)node[right]{$x$};\draw[axis](0,0)--(0,3)node[above]{$y$};
\draw[world,blue]plot[domain=0:5.6,samples=70](\x,{1.4*\x-.25*(\x)^2});
\foreach \x/\lab in {1/t_1,2/t_2,3/t_3,4/t_4,5/t_5}{
\fill(\x,{1.4*\x-.25*(\x)^2})circle(1.5pt);\draw[guide](\x,0)--(\x,{1.4*\x-.25*(\x)^2});\node[below]at(\x,0){$\lab$};}
\node at(3,-.75){相邻标记的时间间隔相等，水平间距相等};
\node[above]at(2.8,2.35){竖直间距和方向随时间改变};
\end{tikzpicture}}
图\ref{k4deep:time}的每个黑点都由同一个时刻的 $x(t_i)$ 和 $y(t_i)$ 配成。相邻黑点的水平位移为 $v_{0x}\Delta t$，所以等时标记在横轴上等距；竖直位移为 $v_y(t_i)\Delta t-\tfrac12g\Delta t^2$，接近最高点时变小，过顶后转为下降。先用落地或碰壁条件求出终止时刻，再截取这些点组成的实际飞行段。

轨迹方程只给空间中的曲线，还须由初始位置、时间方向及边界筛选实际经过的弧段。反过来，只知道轨迹，也不能唯一确定速率：同一空间曲线可用不同时间规律通过。若额外给一个坐标与时间的关系，应先代入轨迹，选定运动分支，再求另一个坐标的导数。

\subsection{速度方向与最高点}
速度为 $(v_0\cos\alpha,v_0\sin\alpha-gt)$，加速度始终竖直向下。最高点只是竖直速度为零，水平速度通常仍非零，因而加速度也不为零。图\ref{k4t-path}中的切线随位置改变，但红色重力箭头方向不变。
\fig{k4t-path}{轨迹切线与恒定重力加速度}{\begin{tikzpicture}\draw[axis](0,0)--(6,0)node[right]{$x$};\draw[axis](0,0)--(0,2.8)node[above]{$y$};\draw[world]plot[domain=0:5.5,samples=60]({\x},{1.1*\x-.2*\x*\x});\draw[flow,blue](1, .9)--(2,1.6)node[above]{$\bm v$};\draw[flow,red](2.75,1.5125)--(2.75,.5)node[right]{$\bm g$};\draw[flow,blue](2.75,1.5125)--(4,1.5125);\node[above]at(2.75,1.7){最高点};\end{tikzpicture}}

用动能关系或两个速度分量，可得 $v^2=v_0^2-2g(y-y_0)$。它能检查同一高度处速率相同，但不能说明速度方向相同。不同高度或碰撞改变运动方向后，应分别建立对应阶段的速度矢量。

\subsection{同高与异高的飞行时间}
发射点高度为 $y_0$、终止水平面为 $y_b$时，飞行时间由 $y_b-y_0=v_{0y}t-gt^2/2$决定。若出发点高于终止面，正根为 $t=[v_{0y}+\sqrt{v_{0y}^2+2g(y_0-y_b)}]/g$；另一个根通常对应过去的延长轨迹。若目标高于发射点，根号非负还只是高度可达的必要条件，同一高度可能在上升、下降两个时刻经过。

同高起落且向上抛射时，$T=2v_0\sin\alpha/g$、水平射程 $S=v_0^2\sin2\alpha/g$。发射角互余的两种射法可有相同射程，却有不同飞行时间和顶点高度。水平抛射从高台出发时，$T=\sqrt{2H/g}$仅由高度决定；如果抛出点与落点不等高，就不能再用同高射程公式。

\subsection{一般恒加速度为何仍产生抛物线}
只要加速度为非零恒矢量，可选一轴沿加速度、另一轴垂直它。垂直方向做匀速运动，平行方向做匀加速运动，消去共同时间即得抛物线。若初速度恰与加速度平行，则垂直方向速度为零，轨迹退化为直线。因而“加速度恒定但不沿初速度”是抛物线轨迹的关键，比限定坐标中的某个方向名称更一般。

已知两个分方向运动时，也要检查所给加速度是否真为恒矢量。若大小恒定但方向始终指向某运动中心，它一般不是抛体运动；把这种加速度代入 $\bm r=\bm r_0+\bm v_0t+\bm at^2/2$会失去方向随时间变化的信息。

\subsection{同一高度处的两时刻与速度方向}
若目标高度比发射点高 $\Delta y$，竖直方程给出
\[t_\pm=\frac{v_{0y}\pm\sqrt{v_{0y}^2-2g\Delta y}}{g}.\]
根号为正且两根为正时，$t_-$对应上升经过，$t_+$对应下降经过。两时刻相对顶点时刻 $t_m=v_{0y}/g$对称，竖直速度等大反向；水平速度却相同，所以两速度一般既不相等也不完全反向。若目标还规定水平位置，就必须把同一个 $t_\pm$代入水平式，不能只满足高度就宣布击中目标。

\fig{k4deep:sameheight}{同高两点的速度竖直分量反号，水平分量相同}{\begin{tikzpicture}
\draw[axis](0,0)--(6,0)node[right]{$x$};\draw[axis](0,0)--(0,3)node[above]{$y$};
\draw[world]plot[domain=0:5.6,samples=70](\x,{1.4*\x-.25*(\x)^2});
\pgfmathsetmacro{\leftroot}{(1.4-sqrt(.76))/.5}\pgfmathsetmacro{\rightroot}{(1.4+sqrt(.76))/.5}
\draw[guide](0,1.2)--(5.6,1.2);\node[left]at(0,1.2){$\Delta y$};
\foreach \x/\sign/\lab in {\leftroot/1/P_-,\rightroot/-1/P_+}{
\fill(\x,1.2)circle(1.5pt)node[below]{$\lab$};\draw[flow,blue](\x,1.2)--++(.8,{.8*\sign*sqrt(.76)});\draw[guide](\x,1.2)--++(.8,0)--++(0,{.8*\sign*sqrt(.76)});}
\node at(2.8,2.6){$v_{y,-}=-v_{y,+}$，$v_{x,-}=v_{x,+}$};
\end{tikzpicture}}
在图\ref{k4deep:sameheight}中，两个水平投影长度相同，竖直投影等长反向。由 $v_y=v_{0y}-gt_\pm$ 立即得到 $v_{y,\pm}=\mp\sqrt{v_{0y}^2-2g\Delta y}$，所以速率相同。两个点的横坐标分别为 $x_0+v_{0x}t_\pm$，一般并不相同；若题目给定的是一个点目标，必须继续筛选横坐标。

对两个同受重力的抛体，可先相减得相对匀速运动，再用任一物体的绝对位置确定碰撞点。相对轨迹直线只能解决两者能否在同一时刻重合，还需检查它们到那一时刻是否已先碰到地面或障碍物。这种先求相对运动、再核对绝对边界的安排，常比直接联立四个分量方程简洁。

\section{斜面与优化问题}
\subsection{沿面和法向坐标}
为避免把竖直距离与法向距离混用，可直接作一次坐标旋转。沿下斜面取 $s=x\cos\theta-y\sin\theta$，外法向取 $n=x\sin\theta+y\cos\theta$。从斜面原点发射时有
\[s=v_0\cos(\alpha+\theta)t+\tfrac12g\sin\theta\,t^2,\qquad
n=v_0\sin(\alpha+\theta)t-\tfrac12g\cos\theta\,t^2.\]
\fig{k4deep:axes}{把重力与初速度分别投影到沿面和法向坐标}{\begin{tikzpicture}
\draw[world](-2,1)--(3,-1.5);\draw[axis](0,0)--(2.6,-1.3)node[right]{$s$};\draw[axis](0,0)--(1.3,2.6)node[above]{$n$};
\draw[flow,red](0,0)--(0,-2.2)node[below]{$\bm g$};\draw[flow,blue](0,0)--(.88,-.44)node[right]{$g\sin\theta$};
\draw[flow,blue](.88,-.44)--(0,-2.2)node[midway,right]{$-g\cos\theta$};
\draw[flow](0,0)--(2,1.3)node[right]{$\bm v_0$};\draw[guide](0,0)--(2,0);\node at(.8,.22){$\alpha$};\node at(1,-.2){$\theta$};
\end{tikzpicture}}
图\ref{k4deep:axes}的蓝色两箭头首尾相接后恢复竖直重力。初速度与沿面正向的夹角是 $\alpha+\theta$，所以法向初速度是 $v_0\sin(\alpha+\theta)$，而非 $v_0\sin\alpha$。把几何投影确定后，两个运动方程都只是各自方向的一维匀变速方程。


$n=0$确定落面时刻，$\dot n=0$确定离面最远时刻。二者依次为 $T=2v_{0n}/(g\cos\theta)$与 $T/2$，因此最大法向距离为 $v_{0n}^2/(2g\cos\theta)$；这一时刻速度沿斜面，通常已经在竖直方向下降。沿面射程则必须将 $T$代入 $s$，或用 $x(T)/\cos\theta$求得，不能把水平射程直接当作沿面长度。

角度的可行域也来自分量：离面要求 $\sin(\alpha+\theta)>0$，向右运动要求 $\cos\alpha>0$。对于向右下的斜面，可以有负的水平仰角，只要发射方向仍在斜面上方。若不先写这些条件，优化时容易把一个数学极值误判成实际可行射法。



从水平坐标计算落点也完全等价：把轨迹与直线 $y=-x\tan\theta$ 联立，并排除发射时刻的零根。图\ref{k4t-plane}中，竖直高度与到斜面的垂直距离不同，后者应沿蓝色法线计算。问“离斜面最远”时，不应直接套最高点的竖直高度。
\fig{k4t-plane}{到斜面的距离沿法线量取}{\begin{tikzpicture}\draw[world](-.5,.25)--(5,-2.5);\draw[world]plot[domain=0:4,samples=50]({\x},{.8*\x-.3*\x*\x});\fill(2,.4)circle(2pt)node[above]{$P$};\draw[blue,dim](2,.4)--(1.44,-.72)node[midway,left]{$h_\perp$};\draw[guide](2,.4)--(2,-1);\node[right]at(2,-.5){竖直距离};\node at(3,-1.2){$\theta$};\end{tikzpicture}}

\subsection{最优条件需要配合可行域}
对发射角求最值之前，先写清落点所在边界、角度区间和碰撞顺序。同高射程最大时的 $45^\circ$ 并不适用于所有斜面或越障问题。对给定目标点求最小初速，可把 $\tan\alpha$ 视为参数，用二次式判别式非负求可达条件；临界时两条射向合并。

越过一座有宽度的建筑，需要同时检查前后两条屋沿，而非只过一个角点。若屋顶水平、起落点等高，最优轨迹的顶点位于屋顶中点正上方；否则把轨迹向有更大余量的一侧平移，仍可能降低所需初速。对称性帮助选参数，最后仍应代回两个屋沿和地面落点检查。

\subsection{给定初速时，同一目标可能对应两个角度}
设发射点为原点，目标为 $(X,Y)$、$X>0$，给定初速 $v_0$。令 $p=\tan\alpha$，并取水平分速度为正。消去时间后，命中条件为
\[Y=Xp-\frac{gX^2}{2v_0^2}(1+p^2).
\]
记 $A=gX^2/(2v_0^2)>0$，则这是 $Ap^2-Xp+(Y+A)=0$。判别式非负时，角度参数为
\[p_\pm=\frac{X\pm\sqrt{X^2-4A(Y+A)}}{2A}.\]
对应的到达时间分别为 $t_\pm=X\sqrt{1+p_\pm^2}/v_0$。应将每一个根都代回检查正时间、角度范围和障碍条件；“二次方程有两个实根”本身还不保证两个射法都被题目允许。

图\ref{k4k-twoangles}取同一目标与同一初速画出高、低两条轨迹。低角度的飞行时间通常较短，高角度轨迹能升得更高；两者经过目标时的竖直速度方向可能不同。临界初速使两根合并，表示满足目标的角度只有一个，而非表示目标必在轨迹顶点。
\fig{k4k-twoangles}{固定初速时，高低两条轨迹可经过同一目标}{\begin{tikzpicture}
\draw[axis](0,0)--(5.5,0)node[right]{$x$};\draw[axis](0,0)--(0,4)node[above]{$y$};
\draw[blue,world]plot[domain=0:4.6,samples=70]({\x},{.5*\x-.0625*1.25*(\x)^2});
\draw[red,world]plot[domain=0:4.2,samples=70]({\x},{3.5*\x-.0625*13.25*(\x)^2});
\fill(4,.75)circle(1.5pt)node[right]{$P$};\draw[guide](4,.75)--(4,0);
\node[blue]at(2,.3){低射};\node[red]at(1.6,3.1){高射};
\end{tikzpicture}}


\subsection{从判别式理解最小初速}
给定目标 $(x,y)$且 $x\ne0$，令 $z=v_0^2>0$，轨迹方程整理为
\[gx^2p^2-2zxp+(gx^2+2zy)=0,\qquad p=\tan\alpha.\]
它对 $p$有实根的条件是
\[\Delta=(-2zx)^2-4gx^2(gx^2+2zy)
 =4x^2(z^2-2gyz-g^2x^2)\ge0.\]
关于 $z$的二次式两根为 $g(y\pm\sqrt{x^2+y^2})$，负号根小于零，所以正的 $z$必须满足 $z\ge g(y+\sqrt{x^2+y^2})$。初速大时可能有一高一低两条射法；逐渐降低初速，两根靠近，临界时合并，继续降低则无实根。于是 $v_{\min}^2=g(y+\sqrt{x^2+y^2})$，临界角由重根 $p=z/(gx)$确定，即 $\tan\alpha=(y+\sqrt{x^2+y^2})/x$。若 $x<0$，还须选 $\cos\alpha<0$的角度分支，保证飞行时间为正。

公式中的平方根是目标到发射点的距离。目标在同一高度时恢复 $v_{\min}^2=g|x|$；若目标几乎在正上方，最优方向趋近竖直。该结果针对自由选择方向且中途无障碍的质点。如果途中有屋顶、墙体或规定发射角区间，达到目标的临界轨迹还可能不符合额外条件，必须继续检查。

\subsection{固定初速时的可达高度与两种射法}
在 $x>0$处，令 $K=gx^2/(2v_0^2)$，轨迹高度是 $y=xp-K(1+p^2)$，其中 $p=\tan\alpha$。完成平方得到
\[y=\frac{v_0^2}{2g}-\frac{gx^2}{2v_0^2}
-K\left(p-\frac{v_0^2}{gx}\right)^2.
\]
平方项说明，无论发射角怎样选，该水平位置的高度都不超过前两项之和。这条上界曲线是各抛物线的包络，并不是某一次固定角度发射的实际轨迹。一个抛体在沿程不同位置不可能不断改变初始角度来始终沿包络运动。

若目标低于上界，角参数通常有两个值，分别对应较低与较高的射法。两种射法的初速大小相同，飞行时间却因水平分量不同而改变。若目标恰在上界上，两根合并，只有一个临界方向。若还设有地面和障碍，应将每个方向恢复为时间方程，检查它是否在首次碰撞前经过目标。仅有空间方程的实根，并不保证实际路径一直有效。

\fig{fv:k4:envelope}{固定初速的一族轨迹共同受包络线限制}{\begin{tikzpicture}[x=1.5cm,y=1.2cm]
\draw[->](0,0)--(3.1,0)node[right]{$gx/v_0^2$};\draw[->](0,0)--(0,1.5)node[above]{$gy/v_0^2$};
\foreach \p in {.5,1,1.7,2.5}{\draw[blue!65,domain=0:{2*\p/(1+\p*\p)},samples=50]plot(\x,{\p*\x-(1+\p*\p)*\x*\x/2});}
\draw[red,thick,domain=0:1.7,samples=60]plot(\x,{.5-.5*\x*\x});\node[red]at(1.5,1.05){包络线};
\end{tikzpicture}}
图\ref{fv:k4:envelope}中的蓝线对应不同发射角，红线给同一横坐标可达到的最大高度。配方法中被减去的平方项表示具体射法相对上界的差；令其为零才找到相切的发射角。因此目标在红线下方只说明有代数候选角，还须核对发射方向、正的飞行时间及中途障碍。若发射角被限制在一个区间，上界也应在该区间内重新求，不能继续沿用全部角度的包络。


\subsection{抛体在最高点附近的局部弯曲}
最高点处 $v_y=0$，速度水平，重力正沿曲线凹侧法向，因此该点曲率半径为 $\rho=v_x^2/g$。这一结果既可从第五讲的自然坐标得到，也可把顶点附近轨迹写成 $y=H-g(x-x_m)^2/(2v_x^2)$读出二阶弯曲。初速度竖直分量主要抬高顶点，水平分量则控制顶点附近的平缓程度，两者在越障优化中承担不同作用。

\fig{k4deep:width}{顶点同高时，水平速度决定轨迹能覆盖多宽的屋顶}{\begin{tikzpicture}
\fill[gray!15](-2.2,0)rectangle(2.2,1);\draw[world](-2.2,0)--(-2.2,1)--(2.2,1)--(2.2,0);
\draw[blue,world]plot[domain=-3:3,samples=70](\x,{2-.18*(\x)^2});
\draw[red,world]plot[domain=-1.9:1.9,samples=60](\x,{2-.5*(\x)^2});
\fill(0,2)circle(1.5pt)node[above]{$M$};\draw[guide](0,0)--(0,2);\node[blue,right]at(2.4,1.65){较大 $v_x$};\node[red,right]at(1.2,.3){较小 $v_x$};
\node[below]at(0,0){同一顶点高度 $H$};
\end{tikzpicture}}
图\ref{k4deep:width}两条曲线的顶点同为 $M$，但红线在到达远侧屋沿以前已经降到屋顶高度以下。由 $y=H-g(x-x_m)^2/(2v_x^2)$，要求屋沿距顶点的水平距离为 $b$ 时仍有 $y\ge h$，就必须满足 $v_x^2\ge gb^2/[2(H-h)]$。这个不等式把“轨迹不能过窄”转成可计算的条件。

因此提高抛物线顶点并不必然降低越过宽屋顶所需的初速；过小的水平分量会使轨迹弯曲过急，远侧屋沿反而难以越过。用顶点高度和水平速度作为参数，能把“足够高”与“足够宽”这两个限制同时写出，而不只凭最高点高于屋顶判断可行。

\subsection{用飞行时间直接求目标所需初速}
对水平距离 $x>0$、高度差 $y$的目标，暂将飞行时间 $t>0$作为未知量。为在该时刻经过目标，初速度分量必须满足
\[v_{0x}=\frac{x}{t},\qquad v_{0y}=\frac yt+\frac{gt}{2}.
\]
于是 $v_0^2=(x^2+y^2)/t^2+gy+g^2t^2/4$。记目标距离 $r=\sqrt{x^2+y^2}$，由两正项乘积固定，可得
\[\frac{r^2}{t^2}+\frac{g^2t^2}{4}\ge gr,\qquad
v_0^2\ge g(r+y).
\]
等号在 $t^2=2r/g$实现，进一步 $\tan\alpha=v_{0y}/v_{0x}=(r+y)/x$。这个推导既给出任何射法都满足的下界，也给出达到下界的时间和方向，与角度判别式互相验证。目标点所需最小初速并不由“瞄准目标直线”给出，因为飞行中还会下落。

在实现最小初速的时刻，$gt^2/2=r$。到达时的竖直速度则是
\[v_y(t)=v_{0y}-gt=\frac yt-\frac{gt}{2}=\frac{y-r}{t}.\]
因为 $x>0$使 $r>|y|$，该值严格为负。因而任意右方有限目标的最小初速射法，经过目标时都已下降；这与判别式重根完全相容。绘制临界轨迹时，顶点应位于目标的左上方，与这一速度方向一致。


\subsection{越过一段障碍需要检查整个区间}
设可飞行范围内的水平屋顶占 $a\le x\le b$，高度为 $h$。轨迹 $y(x)$是开口向下的抛物线，所以在闭区间上的最小值必位于某个端点。证明可直接用凹性：内部点的高度不低于两端高度的线性插值。因此，只要 $y(a)\ge h$且 $y(b)\ge h$，整个水平屋顶便不会比轨迹更高；这解释了为何有限宽建筑题应同时列两个屋沿条件。

对其他形状的障碍，应比较间隙函数 $f(x)=y(x)-y_{\text{障}}(x)$。若屋顶为直斜线，减去一次函数后间隙仍为凹二次式，仍只需检查该段两端；分段直线屋顶则需逐段检查所有接点。若屋顶曲率变化，间隙可能不再凹，内部驻点也可能成为最小间隙。图\ref{k4t-obstacle-coach}中，红色竖线表示两端余量，而蓝色曲线在整个水平屋顶上方。若轨迹被更早的墙或地面截断，后续数学延长线即使满足屋沿不等式，也不属于可行过程。
\fig{k4t-obstacle-coach}{水平屋顶的两端条件与区间内的净空}{\begin{tikzpicture}
\draw[axis](-.2,0)--(6.3,0)node[right]{$x$};
\fill[gray!15](1.4,0)rectangle(4.6,1.4);\draw[world](1.4,0)--(1.4,1.4)--(4.6,1.4)--(4.6,0);
\draw[blue,world]plot[domain=0:6,samples=70]({\x},{.3*\x*(6-\x)});
\draw[red,dim](1.4,1.4)--(1.4,1.932);\draw[red,dim](4.6,1.4)--(4.6,1.932);
\node[below]at(1.4,0){$a$};\node[below]at(4.6,0){$b$};\node at(3,.75){水平屋顶};
\end{tikzpicture}}

以水平宽屋顶为例，写轨迹 $y=H_m-(x-s)^2/(2p)$。两端的高度条件给出 $H_m$的下界，代回 $v_0^2/g=p+2H_m$后才能对 $p,s$共同优化。若只是先猜轨迹对称再求一个候选值，没有排除非对称轨迹可能更优；全局下界与可实现的等号构形应同时给出。

图\ref{k4k-clearance}画出净空函数 $q(x)=y(x)-H$。两端非负与内部开口向下共同保证全程非负；虚线是两端连线，整条曲线都在它上方。若题目要求严格不接触，应把等号的掠过情形说明为极限，而不是自动当成实际通过。
\fig{k4k-clearance}{水平屋顶上方的凹净空函数由两端控制}{\begin{tikzpicture}
\draw[axis](0,0)--(5.3,0)node[right]{$x$};\draw[axis](0,0)--(0,2.8)node[above]{$q=y-H$};
\draw[world]plot[domain=.6:4.5,samples=60]({\x},{.3+.42*(\x-.6)*(4.5-\x)});
\draw[guide](.6,.3)--(4.5,.3);\draw[guide](.6,0)--(.6,.3) (4.5,0)--(4.5,.3);
\node[below]at(.6,0){$a$};\node[below]at(4.5,0){$b$};
\end{tikzpicture}}


\subsection{命中目标与越过墙顶的边界区别}
命中一个点要求轨迹取等号 $y(X)=Y$；越过同位置墙顶若允许掠过，要求 $y(X)\ge Y$，两者在最优边界上给出相同的临界初速。若要求严格正的净空 $y(X)>Y$，则临界速度只是不取到的下确界，实际初速必须更大。若球有非零半径，或墙有厚度，还应把球心的可行域按几何尺寸调整，再建立新的高度或区间条件。

最优解中“等号能否实现”决定最终答案应称最小值还是下确界。两墙周期碰撞的墙角极限、有限屋顶的掠沿构形、给定目标的判别式重根，都是同一边界思想在不同模型中的表现。求出数值以后，应回到具体轨迹核验接触、方向及碰撞顺序。

\section{反弹的分段与展开}
\subsection{光滑完全弹性碰撞改变哪个分量}
碰到固定光滑水平地面时，水平速度不变，竖直速度反号；碰到固定竖直墙时，水平速度反号，竖直速度不变。完全弹性保证速率不变，不代表速度矢量不变。碰后用碰撞点位置和碰后速度作为新初值，才能得到后续轨迹。

\fig{k4deep:reflect}{固定光滑边界只改变速度的法向分量}{\begin{tikzpicture}
\begin{scope}
\draw[world](-1.8,0)--(1.8,0);\draw[flow,blue](-1.2,1.6)--(0,0)node[midway,left]{$\bm v^-$};\draw[flow,red](0,0)--(1.2,1.6)node[midway,right]{$\bm v^+$};
\draw[guide](0,0)--(0,2.3);\node at(0,-.55){水平面：$(v_x,v_y)\to(v_x,-v_y)$};
\end{scope}
\begin{scope}[xshift=6.5cm]
\draw[world](0,-1.8)--(0,1.8);\draw[flow,blue](-1.6,1.2)--(0,0)node[midway,above]{$\bm v^-$};\draw[flow,red](0,0)--(-1.6,-1.2)node[midway,below]{$\bm v^+$};
\draw[guide](-2.2,0)--(.7,0);\node at(0,-2.35){竖直墙：$(v_x,v_y)\to(-v_x,v_y)$};
\end{scope}
\end{tikzpicture}}
图\ref{k4deep:reflect}把碰撞前后速度放在同一个接触点比较。水平面的切线是水平方向，竖直墙的切线是竖直方向；沿这些切线的分量均不变。法向分量反号后，必须从该接触点重新计时，重力依旧竖直向下。

水平抛出的球从高度 $H$ 首次落地需 $T=\sqrt{2H/g}$。反弹到原高度再落地需 $2T$，因此若连续在水平地面反弹，各次触地时刻为 $T,3T,5T,\ldots$。这一规律依赖每次碰撞都无能量损失，也依赖同一水平地面；斜面碰撞通常会改变两个地面坐标分量。

\subsection{反弹后要重新定义局部时间与初速度}
光滑地面上一次完全弹性反弹保持水平速度，把竖直速度反号。若首次触地时刻为 $t_1$、触地点为 $x_1$、触地前竖直速率为 $W$，反弹后的局部时间取 $\tau=t-t_1$，则
\[x=x_1+v_x\tau,\qquad y=W\tau-\frac12g\tau^2.
\]
下一次触地由 $y=0$的正根 $\tau=2W/g$确定。若碰到竖直墙，则竖直分量不跳变，只有 $v_x$反号。重力仍始终竖直向下，不能把整条加速度也镜面对称后继续当作真实空间的重力。


\subsection{水平反弹的统一分段函数}
以水平抛射初始高度 $H$为例，令 $T=\sqrt{2H/g}$。第 $k$次反弹后的竖直段可围绕其顶点时刻 $2kT$写成
\[y(t)=H-\frac g2(t-2kT)^2.
\]
其中 $k=0$只取 $0\le t\le T$；$k\ge1$取 $(2k-1)T\le t\le(2k+1)T$。在连接时刻，左右高度都为零，竖直速度则从 $-gT$跳到 $+gT$。若水平分量恒为 $u$，将 $t=x/u$代入即可得到每一段的空间轨迹和对应水平区间。

图\ref{k4t-bounce-coach}中，$T,3T,5T$为触地时刻，$2T,4T$为反弹后的最高点。两种射法比较墙高时，必须分别判断墙址属于各自哪个 $k$，再联立高度函数。一个根代回高度正确但落在所取分段之外，仍是无效根；这也是反弹问题不能只画一条无限延伸抛物线的原因。
\fig{k4t-bounce-coach}{反弹高度函数的分段与顶点时刻}{\begin{tikzpicture}[x=.85cm,y=.7cm]
\draw[axis](0,0)--(6,0)node[right]{$t/T$};\draw[axis](0,0)--(0,2.6)node[above]{$y$};
\draw[world]plot[domain=0:1,samples=25]({\x},{2*(1-(\x)^2)});
\foreach \k in {2,4}\draw[world]plot[domain={\k-1}:{\k+1},samples=35]({\x},{2*(1-(\x-\k)^2)});
\foreach \x in {1,2,3,4,5}\node[below]at(\x,0){$\x$};\node[left]at(0,2){$H$};
\end{tikzpicture}}

\subsection{展开图与交点筛选}
将地面下方的抛物线关于地面反射，可把一次反弹表示成两段容易比较的轨迹。比较两种射法共同越过的小墙时，应先明确每一条球路在墙处属于第几段，再联立高度。代数求出的交点若落在另一碰撞阶段，应剔除，发射点和共同落点也不一定是所求墙的位置。

周期碰壁问题则要同时满足水平往返周期和竖直反弹周期相等。周期条件只是必要的几何时间约束，还要确认碰壁、触地的先后顺序，以及是否恰好撞到墙角。若题目要求碰撞点处于边界内部，出现在墙角的极限轨迹只能给下确界，不能称为已经达到的最小值。
\subsection{碰撞展开法的适用边界}
反射轨迹与反射速度是同一几何思想的两种表达。固定竖直墙可用镜像空间展开，水平速度在展开图中不反向；但重力竖直方向在各水平镜像中保持一致。对水平地面作镜像时，必须同时正确处理每段抛物线的竖直相位，不能把有重力的整条曲线简单当作无折线的匀速光线。

若碰撞恢复系数小于1，法向碰后速率减小，后续飞行时间不再是固定的 $2T$，连续反弹的高度与时间会形成另一组递推。若地面或墙在运动，反射应先在边界的瞬时参考系中对法向相对速度作反号，再变回地面系。完全弹性描述的是相对固定光滑面的理想规则，不是任意情况下“对地速率永不改变”。

求完后可按实际事件时间排列一张简短的过程式：发射、首次触地、墙顶位置、后续触地、终点。把所求墙的位置代入水平运动求时间，再放回对应竖直段验证高度，能同时检查代数根、碰撞次数和图形先后关系，避免解答只给出一个孤立交点。

